\documentclass[11pt]{article}
\usepackage[margin=0.75in]{geometry}
\usepackage{fontspec}
\usepackage{setspace}
\usepackage{hyperref}
\hypersetup{colorlinks=true, urlcolor=blue, linkcolor=black}
\pagestyle{plain}
\setlength{\parskip}{0.35em}
\setlength{\parindent}{0pt}
\setstretch{0.96}
\setmainfont{Times New Roman}
\usepackage{titlesec}
\titlespacing*{\section}{0pt}{0.52em}{0.2em}

\begin{document}

\begin{center}
{\Large Responses to the Key Selection Criteria}\\[0.3em]
Ke Cheng HSU (Clio)\\[0.15em]
Lecturer in Economics (Level B), Deakin Business School\\
{\small \href{mailto:khsu@connect.ust.hk}{khsu@connect.ust.hk} $\cdot$ \href{https://cliohsu1997.github.io/}{cliohsu1997.github.io}}
\end{center}

\vspace{0.15em}
\section*{1. Ph.D.\ in Economics or equivalent relevant qualifications and experience}

I am a Ph.D.\ candidate in Economics at the Hong Kong University of Science and Technology (HKUST), with expected graduation in June 2027. I also hold an M.Phil.\ in Economics from HKUST (2020--2022) and a B.Sc.\ in Mathematics and Economics with a minor in Actuarial Mathematics from HKUST (2015--2019), where I received the Academic Achievement Medal (awarded to the top 1\% of graduates, 2019) and a University Scholarship (730K HKD, 2015--2019). My thesis supervision committee comprises Associate Professor Pak Hung Au (HKUST), Associate Professor Kohei Kawaguchi (HKUST), and Professor Kei Kawai (University of Tokyo). My doctoral training in microeconomic theory, structural econometrics, and computation (R, Python, C++, Git, and AWS) is complemented by my experience as a Research Assistant for Professor Benjamin Bernard at National Taiwan University (2020), a Teaching Assistant in the Department of Economics at HKUST (2020--2025), Teaching Assistant Coordinator for the HKUST Business School (2025--2026), and Lecturer for the MSc in Economics Math Camp at HKUST (2026).

\section*{2. Field preference, and a developing research record through publications or external funding}

The search is open to all fields and gives preference to macroeconomics, international trade, or microeconomic theory. My research fields are Empirical Industrial Organization and Digital Economics. Industrial organisation is one of the Economics Discipline's focus areas, and my research agenda builds microeconomic and industrial-organization theory models of decisions under uncertainty---where agents make an initial commitment anticipating costly later adjustment---and connects those models to real markets using large transaction and retail datasets.

My theoretical working paper, ``Product Design with Heterogeneous Search Costs'' (with Pak Hung Au; presented at APIOC 2023 and AMES 2023 at NTU Singapore), connects directly to microeconomic theory. Building on Wolinsky's (1986) sequential search model with differentiated products, we introduce search-cost heterogeneity to analyze how firms use observable product design to shape consumer search. In equilibrium, low-search-cost consumers seek niche products and high-search-cost consumers seek standardized products, and a reduction in search costs can raise prices for both designs when this composition shift dominates the search-encouraging effect.

My job market paper, ``Designing Return Policy in E-Commerce: A Trade-off between Purchase and Return Efficiency'' (presented at the HKUST Business School PhD Conference, the CEP Brownbag, and the 2026 China Economics Summer Forum; scheduled for ESAM 2026 and APIOC 2026), studies how return policies balance ex-ante purchase efficiency and ex-post return efficiency. A stricter policy curbs inefficient returns after delivery, but the lower return rate makes a higher price more attractive for the seller and can exclude efficient purchases. Using transaction data from a Taobao apparel retailer across six categories, I estimate a structural model of purchase, return, and return-shipping insurance under monopoly pricing. Automatic return-shipping coverage performs better when margins, net scrap value, and uncertainty are high, whereas optional coverage performs better when those attributes are low; a menu pairing a discounted price with a medium restocking fee balances both margins and outperforms both insurance policies for most categories.

Two further papers are in progress. ``Too Much Unsold Stock?: Environmental Consequences of Price Commitment'' (with Kohei Kawaguchi and Kei Kawai) builds and estimates a dynamic oligopoly model with demand uncertainty on monthly Japanese bookstore data to measure how much unsold stock is returned and discarded under resale price maintenance and how a deposit-refund tax can reduce waste. ``Pass-Through and Incidence with Product Returns'' (with Ce Bian and Liang Zhao) shows that when a higher price raises consumers' return probability, a sales tax or platform commission fee can create negative pass-through, lower the equilibrium price, and raise social welfare by reducing the return rate. As a graduating Ph.D.\ candidate, I am preparing these working papers for journal submission and have been supported to date by competitive institutional funding at HKUST: the Postgraduate Studentship (2020--present), the Red Bird PhD Award (2022), and Research Travel Grants (2023 and 2026).

\section*{3. Teaching at undergraduate and postgraduate level, curriculum development, and Honours or HDR supervision}

From 2020 to 2025 I served as a Teaching Assistant at HKUST across five undergraduate and postgraduate appointments, achieving an overall average evaluation of 4.72 out of 5. At the postgraduate level, I served twice as Teaching Assistant for Ph.D.\ Microeconomic Theory I (ECON5210), receiving 4.85/5 in Fall 2023--24 and 4.82/5 in Fall 2024--25, including 4.85/5 on tutorial materials and activities helping students achieve learning outcomes, 4.83/5 on clarity of explanation, and 4.85/5 on helpful and timely feedback. At the undergraduate level, I assisted in Economics of Uncertainty and Information (ECON4134, 4.83/5; ECON4999T, 4.56/5) and Macroeconomics (ECON2123, 4.53/5). In classroom and office-hour teaching, I begin from concrete economic intuition---such as restaurant queues to motivate herding under imperfect information or consumption vouchers versus cash transfers to motivate fiscal multipliers---before introducing formal equations, and in doctoral tutorials I approach proofs in multiple ways until a student's reasoning is clear.

In 2026 I served as Lecturer for the five-day MSc in Economics Math Camp at HKUST, where I designed the curriculum, lecture notes, and exercises from scratch for incoming postgraduate students from non-mathematics backgrounds. I built the mathematical tools step by step across optimization, linear algebra, calculus, real analysis, and dynamic programming, receiving an average of 4.53/5 across instructor evaluation items (including 4.64/5 on preparation and responsiveness and 4.55/5 on creating a supportive learning environment). Alongside my appointment as Teaching Assistant Coordinator for the HKUST Business School (2025--2026), where I mentored postgraduate students on designing teaching materials and evaluated their mock teaching, my doctoral tutorial teaching and research program in theoretical and empirical industrial organization prepare me to supervise Honours and HDR students at Deakin. I am prepared to teach microeconomic theory, industrial organization, applied microeconomics, information economics, and econometrics from the undergraduate to the Ph.D.\ level, mathematics for economists through the master's level, and macroeconomics and computational methods at the undergraduate level.

\section*{4. Research leadership, administration, and community engagement through impactful research}

My current papers and five-year research agenda---extending return-policy design to joint consumer search and return, product substitutes and complements (such as size and color bracketing in online retail), and two-sided platform competition---address policy and sustainability questions relevant to retail markets, digital platforms, and regulators, and provide a clear foundation for competitive grant applications and industry engagement at Deakin. In particular, the Japanese bookstore project quantifies the environmental waste of unsold inventory discarded under price commitment and evaluates a deposit-refund tax, while the pass-through project analyzes how platform fees and sales taxes affect equilibrium prices, return rates, and consumer welfare. In academic administration and teaching coordination, I served as Teaching Assistant Coordinator for the HKUST Business School (2025--2026), organized and delivered the 2026 MSc Math Camp, and served as an effective link between students and faculty as a teaching assistant (rated 4.48/5 on bridging with the professor and 4.80/5 on proactive reminders of key dates and activities).

\section*{5. Interpersonal skills and working relationships with colleagues}

I have built productive working relationships with faculty co-authors and colleagues across HKUST and the University of Tokyo (Pak Hung Au, Kohei Kawaguchi, Kei Kawai, Ce Bian, and Liang Zhao) and previously worked as a Research Assistant for Benjamin Bernard at National Taiwan University. In departmental service, mentoring postgraduate students on lesson design and mock teaching as Teaching Assistant Coordinator required clear, constructive peer communication. Across my teaching appointments, student evaluations consistently reflect strong interpersonal skills, including 4.76/5 on approachability and friendliness, 4.76/5 on accurate and timely information, and 4.67/5 on responsiveness in undergraduate courses, alongside 4.83/5 on encouraging active participation and 4.82/5 on informative and constructive answers in Ph.D.\ Microeconomic Theory I.

\end{document}
